\thispagestyle{empty}
\chapter*{Abstract}

The term transients encompasses some of the most energetic and exotic objects in the universe. Radio astronomy gives a unique vantage point into many of these transients. Whether it be neutron stars that form in the wake of dead stars, fast, enigmatic bursts at extragalactic distances or chasing potential technological signatures emitted from civilisations on distant worlds. This thesis focuses on searching for transients, both natural and artificial. 

This thesis first details the search for coherent radio emission from a rapidly rotating neutron star (Pulsar) spinning on millisecond timescales in a binary system with a low-mass companion known as a redback. The search candidate redback J0523.5$-$2529 is presented using 34.5 hours of observations with the Parkes and Green Bank Telescopes. No periodic or single-pulse emission is detected, placing stringent limits on the radio luminosity of the system and suggesting either prolonged accretion, unfavourable beaming geometry, or an intrinsically radio-faint neutron star.

Secondly, the largest collection of giant pulses from the Crab to date is outlined. Over 6 years of data from the Irish LOw Frequency Array (I-LOFAR) is analysed to find over 1,000,000 regular pulses and more than 500,000. Per-epoch measurements of dispersion measure and scattering provide new constraints on the temporal evolution of the Crab Nebula and establish a benchmark dataset for future studies of pulsar propagation effects, creating a valuable asset to the community. 

The latter half of this thesis details a series of efforts to search for extraterrestrial technosignatures at low frequencies.  The results from the largest Galactic narrowband technosignature survey are detailed. Using two international stations of the LOw Frequency ARray (LOFAR), a dual-site technosignature survey was carried out, surveying over 1,600,000 targets from \textit{Gaia}. This survey also demonstrates the effectiveness of using dual-site observations for discriminating against radio frequency interference and establishes a scalable framework for future low‑frequency technosignature programmes.

Expanding on the first survey, the development and initial results of the LOw Frequency pulsar, FRB and Technosignature Survey (LOFTS) are described. LOFTS takes lessons learned from the survey and increases both the analysis and observation time. At the time of writing, the current analysed sample of  \textit{Gaia} targets is over 3,000,000. The survey also includes the search for natural transients such as fast radio bursts, flare stars and pulsars. LOFTS also expands the pipelines used in conventional narrowband signature searches and incorporates methods to account for pulse broadening seen at low frequency due to scintillation, along with the development of a machine learning architecture for narrowband signal detection. 

Finally, work describing the development of the Lunar Farside Technology and Transients telescope is described. Here, the feasibility of exploring an interference-free and unique frequency parameter space for novel transients is described. 

Collectively, this thesis demonstrates how low-frequency astronomy can expand the discovery space for both natural and artificial radio transients and help develop frameworks for future low-frequency transient science and surveys from both the ground and the farside of the Moon.  